% 原书第 7—13 页；第 1.4 节及两个小节、图 1.3—1.4、表 1.1。
\FloatBarrier
\section{两个例子}
\label{sec:two-examples}

\subsection{模式识别}
\label{sec:pattern-recognition}

第一个例子是\textbf{光学字符识别（optical character recognition）}，这是机器学习中已经得到充分研究的问题。它并不是介绍因果结构时通常会采用的例子，却可能有助于熟悉机器学习的读者理解相关概念。下面我们将描述两个因果模型，它们都能在两个随机变量之间产生相依关系。在这个例子中，我们把这两个变量分别记为手写数字图像 $X$ 和\textbf{类别标签（class label）} $Y$。两个模型的底层因果结构不同，却会导出相同的统计结构。

\textbf{模型（i）}假设，我们通过如下过程生成每一对观测值：向书写者依次提供类别标签 $y$，并要求其每次都画出与标签对应的手写数字图像 $x$。我们假定书写者会尽力正确完成任务，但其感知类别标签的过程，以及执行绘图动作的过程，都可能存在噪声。我们可以把图像 $X$ 写成类别标签 $Y$ 和某个独立噪声 $N_X$ 的函数 $f$，并将这一函数理解为生成图像的\textbf{机制（mechanism）}。这里，$Y$ 也被视为随机变量，见图 1.3 左图。于是，我们可以根据 $P_Y$、$P_{N_X}$ 和 $f$ 计算 $P_{X,Y}$。这一分布称为\textbf{观测分布（observational distribution）}；“观测”意味着，我们只是被动地观察系统，而没有进行干预。$X$ 和 $Y$ 是相依的随机变量，因此，我们能够从观测中学习由 $x$ 到 $y$ 的映射，并且根据图像 $x$ 预测正确标签 $y$，其准确性优于随机猜测。

在这一因果结构中，有两类可能的干预，分别会导出相应的\textbf{干预分布（intervention distribution）}。\footnote{第 6.3 节将说明，更一般地说，干预可以理解为对函数和随机变量的改变。}
如果干预生成后的图像 $X$，例如修改图像，或者在图像生成后将其替换为另一幅图像，那么，先前提供给书写者并记录在数据集中的类别标签不会受到影响。形式上，改变 $X$ 不会影响 $Y$，因为 $Y\coloneqq N_Y$。相反，对 $Y$ 进行干预，意味着改变提供给书写者的类别标签，这会显著影响生成的图像。形式上，改变 $Y$ 会影响 $X$，因为 $X\coloneqq f(Y,N_X)$。图中的箭头体现了这种方向性，我们用它表示\textbf{直接因果关系（direct causation）}。

在另一种\textbf{模型（ii）}中，我们不向书写者提供类别标签，而是让书写者自行决定要写哪个数字，并同时记录相应的类别标签。在这种情况下，图像 $X$ 和记录下来的类别标签 $Y$ 都是书写者意图的函数。我们将这一书写意图记为 $Z$，同样把它看作随机变量。为使模型具有一般性，我们不仅假设图像生成过程含有噪声，也假设类别标签的记录过程含有噪声；这些噪声项仍然相互独立，见图 1.3 右图。需要注意，只要适当选择函数和噪声项，就能使这一模型导出的观测分布 $P_{X,Y}$ 与模型（i）的观测分布完全相同。\footnote{事实上，命题 4.1 表明，任意联合分布 $P_{X,Y}$ 都可以由这两个模型分别导出。}

\begin{figure}[htbp]
\centering
\begin{tikzpicture}[causal,x=1cm,y=1cm,font=\small]
 \node[vtx] (f13aX) at (-5.4,0) {$X$};
 \node[vtx] (f13aY) at (-2.2,0) {$Y$};
 \draw[ed] (f13aY)-- node[above=2mm] {$X\coloneqq f(Y,N_X)$} (f13aX);
 \node at (-2.2,-.85) {“2”};
 \node[align=center] at (-3.8,-1.5) {模型（i）；$Y,N_X$ 相互独立};
 \node[vtx] (f13bX) at (1.2,0) {$X$};
 \node[vtx] (f13bY) at (5.2,0) {$Y$};
 \node[latent] (f13bZ) at (3.2,1.8) {$Z$};
 \node[above=2mm] at (f13bZ.north) {书写意图};
 \draw[ed] (f13bZ)-- node[pos=.52,above left=-.5mm] {$X\coloneqq g(Z,M_X)$} (f13bX);
 \draw[ed] (f13bZ)-- node[pos=.52,above right=-.5mm] {$Y\coloneqq h(Z,M_Y)$} (f13bY);
 \node at (5.2,-.85) {“2”};
 \node[align=center] at (3.2,-1.5) {模型（ii）；$Z,M_X,M_Y$ 相互独立};
 % 用矢量笔迹表现原图的手写数字样例。
 \foreach \digitx in {-5.4,1.2}{
   \begin{scope}[shift={(\digitx,-.85)}]
    \draw[line width=1.3pt,line cap=round,line join=round]
      (-.13,.12) .. controls (-.09,.26) and (.14,.24) .. (.16,.10)
      .. controls (.18,.02) and (-.11,-.04) .. (-.13,-.15)
      .. controls (.00,-.12) and (.11,-.18) .. (.18,-.12);
   \end{scope}
 }
\end{tikzpicture}
\caption{手写数字数据集的两种结构因果模型。在左侧模型（i）中，书写者获得类别标签 $Y$，并据此生成图像 $X$；在右侧模型（ii）中，书写者自行决定要写哪一类数字（$Z$），再生成图像和类别标签。选择适当的函数 $f,g,h$ 以及噪声变量 $N_X,M_X,M_Y,Z$，就能使两个模型产生相同的可观测分布 $P_{X,Y}$，但它们在干预下仍有不同表现，见第 1.4.1 节。}
\label{fig:1.3}
\end{figure}


现在考察模型（ii）中可能的干预。如果干预图像 $X$，情况与刚才相同，类别标签 $Y$ 不受影响。但是，如果干预类别标签 $Y$，即修改书写者记录下来的标签，那么与前一种模型不同，这次并不会影响图像。

总之，如果不对所涉及的函数类和分布类施加限制，那么模型（i）与模型（ii）可以在 $X$ 和 $Y$ 上诱导出相同的观测分布，却产生不同的干预分布。在一切都由 $P_{X,Y}$ 导出的纯概率描述中，我们看不到这种差别。然而，只要纳入关于 $P_{X,Y}$ 如何产生的结构性知识，特别是图结构、函数和噪声项，就能够讨论这一差别。

模型（i）和模型（ii）都是\textbf{结构因果模型（structural causal model，SCM）}的例子；这类模型有时也称为\textbf{结构方程模型（structural equation model）}［例如 Aldrich, 1989；Hoover, 2008；Pearl, 2009；Pearl et al., 2016］。在结构因果模型中，所有相依关系都由函数生成：这些函数根据其他变量计算某个变量的值。关键在于，应把这些函数读作\textbf{赋值（assignment）}，也就是按计算机科学中的函数来理解，而不能仅将它们看作数学等式。我们通常把这些函数视为对物理机制的建模。一个结构因果模型会导出所有可观测变量的联合分布。我们已经看到，同一个分布可以由不同的结构因果模型生成，因此，从结构因果模型转向相应的概率模型时，有关干预效应的信息，以及第 6.4 节将讨论的\textbf{反事实（counterfactual）}信息，都可能丢失。本书以结构因果模型为出发点，并尝试在这一基础上展开后续内容。

最后，就这个例子再说明两点。

第一，图 1.3 很好地展示了 Reichenbach 的共同原因原理。$X$ 和 $Y$ 的相依关系可以有多种因果解释。在右图中，给定 $Z$ 后，$X$ 与 $Y$ 就相互独立：图像与标签之间不存在书写意图之外的共享信息。

第二，有时人们认为，只有把时间概念纳入考虑，才能讨论因果性。时间确实在上述例子中发挥了作用，例如，它使我们能够排除“干预 $X$ 会影响类别标签”这种可能性。不过，这并没有什么问题：统计数据集由某个随时间展开的过程生成，本来就是很常见的情形。例如，在模型（i）中，$X$ 与 $Y$ 之间的统计相依关系，源于一个动态过程。书写者先读取标签、规划动作；这涉及大脑中的复杂活动，最终由肌肉和笔执行动作。我们对这一过程的理解还不完整，但它确实是一个随时间展开的物理动态过程，其最终结果是 $X$ 和 $Y$ 之间一个非平凡的联合分布。进行统计学习时，我们只关心这一最终结果。

因此，不仅因果结构，纯粹的概率结构也可以由随时间展开的过程产生；甚至可以认为，这归根结底是它们唯一可能的产生方式。不过，在这两类描述中，暂时不考虑时间往往都有助于分析。在统计学中，讨论统计相依等概念时，通常不必涉及时间；在因果模型中，讨论干预效应时，也往往不必涉及时间。这两个描述层次都可以看作对底层物理模型的抽象；后者更加精确，对现实的描述也比它们更完整，见表 1.1。此外，模型中的变量未必对应明确界定的时点。例如，心理学家研究学生学习动机与学习表现之间的统计关系或因果关系时，就很难为这两个变量分别指定某个具体时点。对应明确时点的测量，反而更常见于物理学、化学等“硬科学”。

\begin{table}[htbp]
\centering
\small
\setlength{\tabcolsep}{4pt}
\begin{tabularx}{\textwidth}{>{\raggedright\arraybackslash}p{3.0cm}*{5}{>{\centering\arraybackslash}X}}
 \toprule
 模型 & 独立同分布\newline 条件下预测 & 分布变化或\newline 干预下预测
      & 回答反事实\newline 问题 & 获得物理\newline 层面的认识 & 从数据\newline 中学习 \\
 \midrule
 机理／物理模型\newline （如第 2.3 节） & 是 & 是 & 是 & 是 & ? \\
 \addlinespace[3pt]
 结构因果模型\newline （如第 6.2 节） & 是 & 是 & 是 & ? & ? \\
 \addlinespace[3pt]
 因果图模型\newline （如第 6.5.2 节） & 是 & 是 & 否 & ? & ? \\
 \addlinespace[3pt]
 统计模型\newline （如第 1.2 节） & 是 & 否 & 否 & 否 & 是 \\
 \bottomrule
\end{tabularx}
\caption{模型的一个简单分类。最详细的模型（最上行）是机理模型或物理模型，通常涉及一组微分方程。另一个极端（最下行）是纯粹的统计模型：这种模型可以从数据中学习，但除了刻画伴随现象之间的关联之外，往往难以提供更多认识。因果模型可以看作介于两者之间的描述：它不再追求完整再现物理细节，却保留了回答某些干预或反事实问题的能力。物理模型与结构因果模型之间的联系，参见 Mooij et al. [2013]；关于干预的讨论，参见第 6.3 节。}
\label{tab:1.1}
\end{table}


\FloatBarrier
\subsection{基因扰动}
\label{sec:gene-perturbation}

第 1.4.1 节已经说明，不同的因果结构会产生不同的干预分布。有时，我们确实希望预测某个随机变量在这类干预下的取值。下面考虑一个来自遗传学的例子，尽管它在某些方面作了过度简化。假设我们已经获得基因 A 的活性数据，以及与之对应的某个\textbf{表型（phenotype）}的测量值。图 1.4 左上图给出了一个示例数据集。可以看到，这两个变量高度相关。这种相关可以用于传统预测：如果观测到基因 A 的活性在 6 左右，我们就会以较高概率预期表型的取值落在 12 至 16 之间。基因 B 的情形也类似，见左下图。

另一方面，我们也可能关心敲除基因 A 后的表型，也就是将该基因的活性设为 0 后会发生什么。\footnote{为简单起见，假设我们能够获得基因的真实活性，不存在测量噪声。}
然而，在完全不了解因果结构的情况下，我们无法给出有实质内容的回答。如果基因 A 对表型具有因果影响，那么干预后，表型预计会发生显著变化，见右上图。事实上，此时甚至可能仍然使用从观测数据中学到的同一个线性模型。反之，如果存在某个共同原因，例如第三个基因 C，同时影响基因 B 的活性和表型，那么对基因 B 的干预就不会影响表型，见右下图。

与模式识别的例子类似，这里的模型也经过适当选择，使基因 A 与表型的联合分布，等于基因 B 与表型的联合分布。因此，仅凭观测数据，无论样本量是否趋于无穷，我们都无法区分图中上下两种情形。归结起来，如果不愿采用因果性方面的概念，那么，对于“敲除基因后，表型会怎样”这一预测问题，我们就只能回答：“不知道。”

\begin{figure}[htbp]
\centering
\definecolor{interventionred}{RGB}{237,28,36}
\begin{tikzpicture}[font=\footnotesize,>={Stealth[length=1.7mm,width=1.1mm]}]
 % 每个面板均使用原 PDF 矢量图中读取的 100 个散点及红色标记坐标。
 \foreach \panel/\letter/\xoffset/\yoffset in {1/A/0/0,2/A/7.5/0,3/B/0/-6.2,4/B/7.5/-6.2}{
  \begin{scope}[shift={(\xoffset,\yoffset)},x=.55cm,y=.23cm]
   \foreach \xtick in {0,2,4,6,8}{
    \draw[gray!25,densely dotted,thin] (\xtick,-.8)--(\xtick,19.8);
   }
   \foreach \ytick in {0,5,10,15}{
    \draw[gray!25,densely dotted,thin] (-1.4,\ytick)--(8.8,\ytick);
   }
   \input{data/figure_1_4_panel_\panel}
   \input{data/figure_1_4_red_panel_\panel}
   \draw[gray,thin] (-1.4,19.8)--(-1.4,-.8)--(8.8,-.8);
   \foreach \xtick in {0,2,4,6,8}{
    \draw[gray,thin] (\xtick,-.8)--(\xtick,-1.45) node[below=2pt,text=black] {$\xtick$};
   }
   \foreach \ytick in {0,5,10,15}{
    \draw[gray,thin] (-1.4,\ytick)--(-1.7,\ytick) node[left=2pt,text=black] {$\ytick$};
   }
   \node at (3.7,-4.15) {基因 \letter\ 的活性};
   \node[rotate=90] at (-3.05,9.5) {表型};
   \ifnum\panel=1
    \node[interventionred,font=\large] at (-.48,5) {?};
   \fi
   \ifnum\panel=3
    \node[interventionred,font=\large] at (-.48,5) {?};
   \fi
   \ifnum\panel=2
    \node[inner sep=1pt] (f14aGene) at (3.1,1.5) {基因 A};
    \node[inner sep=1pt] (f14aPheno) at (7.2,1.5) {表型};
    \draw[->,thin] (f14aGene)--(f14aPheno);
   \fi
   \ifnum\panel=4
    \node[inner sep=1pt] (f14bCommon) at (5.4,5.9) {混杂因子};
    \node[inner sep=1pt] (f14bGene) at (4.0,1.5) {基因 B};
    \node[inner sep=1pt] (f14bPheno) at (7.6,1.5) {表型};
    \draw[->,thin] (f14bCommon)--(f14bGene);
    \draw[->,thin] (f14bCommon)--(f14bPheno);
   \fi
  \end{scope}
 }
\end{tikzpicture}
\caption{两个基因的活性（上排：基因 A；下排：基因 B）都与表型高度相关，黑色散点表示观测数据。然而，敲除基因，即将其活性设为 0 时，对表型的最佳预测（左列）取决于因果结构（右列）。如果基因与表型之间的相关由某个共同原因造成，那么干预后表型预计仍保持通常的表现（右下图）；如果基因对表型具有因果影响，干预就会明显改变表型的取值（右上图）。本图的构思来自 Peters et al. [2016]。}
\label{fig:1.4}
\end{figure}

\FloatBarrier
